% ============================================================
% Informe TP0 - Cuartetos de Números Primos
% Compilar con pdflatex (o XeLaTeX) en Overleaf.
%
% Para que las figuras aparezcan, subí a Overleaf (en la misma
% carpeta que este .tex, o ajustá los nombres):
%   - grafico_busqueda_base.png
%   - grafico_busqueda_optimizada.png
%   - grafico_busqueda_comparacion.png
% ============================================================
\documentclass[12pt,a4paper]{article}

\usepackage[utf8]{inputenc}
\usepackage[T1]{fontenc}
\usepackage[spanish,es-tabla]{babel}
\usepackage{amsmath,amssymb}
\usepackage{graphicx}
\usepackage{booktabs}
\usepackage{geometry}
\usepackage{float}
\usepackage{listings}
\usepackage{caption}
\usepackage[hidelinks]{hyperref}

\geometry{margin=2.5cm}

% Estilo del pseudocódigo
\lstset{
  basicstyle=\ttfamily\small,
  frame=single,
  framerule=0.5pt,
  columns=fullflexible,
  keepspaces=true,
  xleftmargin=1.5em,
}

\title{\textbf{Informe TP0}\\[0.4em] \large Cuartetos de Números Primos}
\author{Francisco Oscar Ballerio}
\date{Agosto 2026}

\begin{document}

\maketitle

\section{Código Inicial}

La solución al ejercicio está compuesta por dos funciones: \texttt{esprimo} y \texttt{prog}.
\begin{itemize}
  \item \texttt{esPrimo} recibe un número entero $n$ y busca la cantidad de divisores de este en el conjunto de números enteros $[2, n-1]$ recorriendo cada elemento. Si esta cantidad es 0 significa que el número es primo y, de tener al menos 1 divisor, el número no es primo.
  \item \texttt{prog} muestra el tiempo que tarda en imprimir por pantalla los cuartetos de números primos encontrados usando un bucle que va desde 11 hasta 1000000 y, con cada iteración, llama a \texttt{esPrimo} para verificar si el número y sus 3 hermanos son primos.
\end{itemize}

\subsection{Supuestos y Condiciones del Algoritmo}

Para que el algoritmo funcione, se necesita:
\begin{itemize}
  \item Que el límite ($n$) sea un entero positivo.
\end{itemize}

El algoritmo:
\begin{itemize}
  \item Busca cuartetos de primos en la misma decena.
  \item Todos los elementos de la solución serán menores que el límite.
  \item Los tiempos de ejecución dependen del hardware y de la cantidad de tareas que el OS está ejecutando al momento de correr el código.
  \item La versión de Python es 3.10.
\end{itemize}

\subsection{Limitaciones en el algoritmo dado}

\textbf{esprimo:}
\begin{itemize}
  \item Recorre todos los números hasta $n-1$.
  \item El bucle \texttt{while} es menos eficiente que un \texttt{for}.
  \item Cuenta la cantidad de divisores que tiene el número.
  \item Compara si la cantidad de divisores es 0 para devolver algo cuando podría devolver directamente la comparación, pues devuelve un \texttt{bool}.
\end{itemize}

\textbf{prog:}
\begin{itemize}
  \item El bucle recorre todos los números de 11 a 1000000. Lo que genera:
  \begin{itemize}
    \item Se saltea el primer cuarteto de primos.
    \item Chequea si un número es primo varias veces.
    \item Sobrepasa la decena. Ej.:
    \begin{verbatim}
i = 11  -> esprimo(11) and esprimo(13) and esprimo(17) and esprimo(19)
i = 13  -> esprimo(13) and esprimo(15) and esprimo(19) and esprimo(21)
    \end{verbatim}
  \end{itemize}
  \item Cada vez que se encuentra un cuarteto, se imprime. Lo que genera operaciones de I/O que ralentizan la búsqueda de los mismos. Imprimir causa que el OS tenga que manejar operaciones de I/O, llamadas al kernel y cambios de contexto para poder hacerlo en cada iteración.
\end{itemize}

\section{Diseño y Mejora del Algoritmo}

\subsection{Optimizando \texttt{esprimo}}

\begin{enumerate}
  \item De bucle \texttt{while} pasaremos a usar un \texttt{for} que directamente ejecuta la iteración en C. \texttt{while} requiere una condición que se verifique en cada iteración con Python.

  \item De la mano del cambio del \texttt{for}, agregamos un \texttt{range(inicio, fin, salto)} que devuelve una lista con el intervalo $[inicio, fin]$.

  \item Buscamos divisores entre 2 y $\sqrt{x}$, pues si $n = a \cdot b$ entonces alguno de los dos es menor o igual a $\sqrt{x}$. En este caso lo que queremos es encontrar al menos 1 divisor para $n$, entonces si en $[1, \sqrt{n}]$ no encontramos ninguno, en $[\sqrt{n}, n-1]$ tampoco habrá: si $a, b$ alguno es mayor a $\sqrt{n}$, entonces necesariamente el otro debe ser menor a $\sqrt{n}$.

  De la mano de este cambio viene el uso de la biblioteca integrada \texttt{math} para calcular la raíz. Esta implementa el cálculo y el casteo a \texttt{int} en C.

  \item Si al menos un número de $[1, \sqrt{n}]$ divide a $n$, entonces no es primo. Por eso, al encontrar un divisor cortamos el bucle.

  \item División optimizada:
  \begin{quote}
    Cualquier número entero mayor a 3 puede representarse algebraicamente bajo una de las formas $6k$, $6k+1$, $6k+2$, $6k+3$, $6k+4$ o $6k+5$ (para $k \ge 1$):
    \begin{itemize}
      \item Las formas $6k$, $6k+2$ y $6k+4$ son pares (divisibles por 2), por lo que son números compuestos.
      \item La forma $6k+3$ es divisible por 3, siendo también compuesta.
      \item Las únicas formas restantes son $6k+1$ y $6k+5$, donde $6k+5$ equivale algebraicamente a $6(k+1)-1$, quedando ambas cubiertas bajo el patrón $6k \pm 1$.
    \end{itemize}
    Dado que los números divisibles por 2 o 3 se pueden descartar en una verificación inicial (en este caso los obviamos por conveniencia del algoritmo), para determinar la primalidad basta con evaluar únicamente los divisores de la forma $6k \pm 1$ hasta el límite $\sqrt{n}$ \cite{geeksforgeeks}.
  \end{quote}

  Esto nos permite optimizar el algoritmo \texttt{esPrimo} y reducir significativamente la cantidad de operaciones necesarias. Finalmente, lo mejor que se pudo mejorar la complejidad de este algoritmo es $O(\sqrt{n})$.
\end{enumerate}

\subsection{Optimizando \texttt{prog}}

\begin{enumerate}
  \item Restricción del espacio de búsqueda para cuartetos de primos:
  \begin{quote}
    Un cuarteto de números primos es un conjunto de la forma $\{p, p+2, p+6, p+8\}$. Todo cuarteto con $p > 5$ cumple con las siguientes propiedades aritméticas ($\{5, 7, 11, 13\}$ es el único que no cumple, pero como no están en la misma decena no nos importa):
    \begin{itemize}
      \item Ninguno de sus elementos puede ser divisible por 2, 3 o 5.
      \item En aritmética modular, esto restringe la posición del primer elemento a una única clase de congruencia módulo 30 ($2 \times 3 \times 5 = 30$), específicamente $p \equiv 11 \pmod{30}$.
    \end{itemize}
    Como consecuencia, la distancia mínima entre el inicio de dos cuartetos consecutivos $\{p, \dots\}$ y $\{q, \dots\}$ es $q - p = 30$, y toda separación entre candidatos válidos debe ser forzosamente un múltiplo de 30 \cite{wikipedia}.
  \end{quote}

  Esta propiedad permite optimizar el algoritmo de búsqueda, descartando la evaluación de valores intermedios y avanzando las iteraciones en saltos de 30 (o evaluando únicamente candidatos $p \equiv 11 \pmod{30}$).

  \item Para solucionar la ralentización de I/O metemos los candidatos en una lista de cuatruplas $(i, i+2, i+6, i+8)$ para imprimirlas todas juntas luego de encontrar \emph{todos} los cuartetos.

  \item El crear la lista de cuartetos \emph{inline} aumenta la performance del código. La implementación de este método está en CPython y evita usar \texttt{.append()} repetidas veces.
\end{enumerate}

\section{Pseudocódigo y Estructuras de Datos}

\begin{itemize}
  \item \textbf{Lista:}
    \begin{itemize}
      \item En la función \texttt{prog}, \texttt{resultados} es una lista con los cuartetos a imprimir.
    \end{itemize}
  \item \textbf{Tupla:}
    \begin{itemize}
      \item Los cuartetos son almacenados como \texttt{tuple[int, int, int, int]} en la lista de resultados.
    \end{itemize}
\end{itemize}

\subsection{\texttt{esprimo}}

\begin{lstlisting}
recibo numero
limite = raiz_cuadrada(numero)
para cada num (desde 5, hasta limite, incrementando 6) {
  si resto(n / num) == 0 o si resto(n / (num + 2)) == 0 {
    devuelvo Falso
  }
}
devuelvo verdadero
\end{lstlisting}

\subsection{\texttt{prog}}

\begin{lstlisting}
resultados = lista[tupla(int, int, int, int)]
para cada i (desde 11, hasta 1000000, incrementando 30) {
  si esprimo(i) y esprimo(i + 2) y esprimo(i + 6) y esprimo(i + 8) {
    lista.agregar(tupla(i, i+2, i+6, i+8))
  }
}
para cada tupla de resultados {
  imprimir tupla
}
imprimir tiempo de ejecucion
\end{lstlisting}

\section{Complejidad Temporal}

\subsection{Original}

\textbf{\texttt{esprimo}}: loopea desde 2 hasta $n-1$, es decir, tiene una complejidad de $O(n)$.

\textbf{\texttt{prog}}: tiene una complejidad de $O(n^2)$, ya que llama a \texttt{esPrimo} para cada número en el rango $[11, 1000000]$.

\subsection{Optimizado}

\textbf{\texttt{esprimo}}: el algoritmo recorre de 6 en 6 los números desde el 5 hasta $\sqrt{n}$. O sea, se realizan una cantidad de $\sqrt{n} / 6$ operaciones, lo que da $O(\sqrt{n} / 6)$ operaciones, que es lo mismo que $O(\sqrt{n})$.

\textbf{\texttt{prog}}: gracias a que \texttt{esprimo} mejoró su complejidad, este algoritmo también lo hizo. Pero la base de la complejidad sigue siendo la misma, pues la cantidad de llamadas a \texttt{esprimo} depende de $n$. Finalmente se realizan $n/30 \cdot 4$ llamadas a \texttt{esprimo}:
\[
O(n/30) \cdot 4 \cdot O(\sqrt{n}) = O(n) \cdot O(\sqrt{n}) = O(n^{3/2})
\]

\section{Comparación de Tiempos de Ejecución}

En la teoría conocemos la complejidad de ambos algoritmos. Ahora queremos comprobar empíricamente que estas complejidades sean correctas. Por lo tanto usaremos el paper proporcionado por la cátedra para aproximar por cuadrados mínimos y verificar que el error sea relativamente bajo.

\subsection{Base}

\begin{table}[H]
\centering
\caption{Tiempos de ejecución medidos para la versión original (base.py).}
\begin{tabular}{@{}rr@{}}
\toprule
$N$ & Tiempo promedio (s) \\
\midrule
499  & 0.0088 \\
626  & 0.0155 \\
785  & 0.0245 \\
984  & 0.0347 \\
1234 & 0.0576 \\
1547 & 0.0928 \\
1939 & 0.1528 \\
2430 & 0.2431 \\
3046 & 0.3582 \\
3818 & 0.5568 \\
4786 & 0.8202 \\
5999 & 1.2352 \\
\bottomrule
\end{tabular}
\end{table}

\begin{table}[H]
\centering
\caption{Ajuste por cuadrados mínimos (base.py).}
\begin{tabular}{@{}lllll@{}}
\toprule
Algoritmo & Modelo & $c_1$ & $c_2$ & Error cuadrático \\
\midrule
busqueda\_base & $c_1 n^2 + c_2$ & $3.47\times10^{-8}$ & $1.48\times10^{-2}$ & $4.12\times10^{-3}$ \\
\bottomrule
\end{tabular}
\end{table}

\begin{figure}[H]
\centering
\includegraphics[width=0.8\linewidth]{grafico_busqueda_base.png}
\caption{Búsqueda de cuartetos (base.py) — ajuste $O(n^2)$.}
\end{figure}

\subsection{Optimizado}

\begin{table}[H]
\centering
\caption{Tiempos de ejecución medidos para la versión optimizada (tp.py).}
\begin{tabular}{@{}rr@{}}
\toprule
$N$ & Tiempo promedio (s) \\
\midrule
100000  & 0.00535 \\
142709  & 0.00939 \\
203659  & 0.01466 \\
290640  & 0.02446 \\
414769  & 0.04034 \\
591914  & 0.06298 \\
844716  & 0.09753 \\
1205487 & 0.16287 \\
1720341 & 0.27145 \\
2455084 & 0.45752 \\
3503629 & 0.74358 \\
4999999 & 1.19270 \\
\bottomrule
\end{tabular}
\end{table}

\begin{table}[H]
\centering
\caption{Ajuste por cuadrados mínimos (tp.py).}
\begin{tabular}{@{}lllll@{}}
\toprule
Algoritmo & Modelo & $c_1$ & $c_2$ & Error cuadrático \\
\midrule
busqueda\_optimizada & $c_1 n^{3/2} + c_2$ & $1.08\times10^{-10}$ & $1.46\times10^{-2}$ & $2.67\times10^{-3}$ \\
\bottomrule
\end{tabular}
\end{table}

\begin{figure}[H]
\centering
\includegraphics[width=0.8\linewidth]{grafico_busqueda_optimizada.png}
\caption{Búsqueda de cuartetos (tp.py) — ajuste $O(n^{1.5})$.}
\end{figure}

\subsection{Conclusión}

\begin{figure}[H]
\centering
\includegraphics[width=0.9\linewidth]{grafico_busqueda_comparacion.png}
\caption{Comparación de tiempos: base.py vs tp.py (escala log-log).}
\end{figure}

El algoritmo optimizado, en el mismo tiempo que el original, puede manejar inputs unas 830 veces más grandes. Siendo que el algoritmo base encuentra los cuartetos menores a 6000 en 1.2\,s y el optimizado encuentra los menores a 6000000 en el mismo tiempo.

\section{Referencias}

\begin{thebibliography}{9}
\bibitem{geeksforgeeks}
GeeksforGeeks. (2026, 25 de agosto). \emph{Program for prime number check}.
\url{https://www.geeksforgeeks.org/dsa/check-for-prime-number/}

\bibitem{wikipedia}
Prime quadruplet. (s.f.). En \emph{Wikipedia}. Recuperado el 26 de agosto de 2026, de
\url{https://en.wikipedia.org/wiki/Prime_quadruplet\#Prime_k-tuples}
\end{thebibliography}

\end{document}
